\section{Conclusioni}

\subsection{Milestone 1}

La Milestone 1 ha prodotto un dataset strutturato di 15.338 snapshot di classi Java per
il progetto Apache OpenJPA, corredato di 19 metriche e
dell'etichetta di bugginess. Il processo di costruzione ha richiesto numerose scelte
metodologiche riguardanti il filtraggio delle versioni, la sanitizzazione dei ticket, la
gestione delle versioni prive di tag Git e la definizione dello shift temporale per i code
smell. Le scelte adottate privilegiano la pulizia e la coerenza dei dati rispetto alla
massimizzazione della dimensione del dataset.

\subsection{Milestone 2}

Il run preliminare (27 combinazioni, SMOTE e wrapper disabilitati) evidenzia che
RandomForest domina su tutte le metriche di ranking (AUC $= 0{,}94$, NPofB20 $= 0{,}88$
senza balancing), confermando le aspettative della letteratura di defect prediction. Le
sue probabilità ben calibrate --- prodotto dell'averaging su 100 alberi --- lo rendono
il classificatore più affidabile per la prioritizzazione delle code review su questo dataset.

La feature selection con soglia $0{,}0$ si è rivelata un no-op: tutte le 19 feature
portano informazione sulla bugginess e nessuna viene eliminata dai metodi filter. I metodi
wrapper potrebbero identificare sottoinsiemi più compatti, ma i tempi di esecuzione ne
hanno finora impedito l'utilizzo.

Il trade-off tra Recall e Kappa introdotto dall'Undersampling (Recall RF: $0{,}52 \to
0{,}86$; Kappa: $0{,}58 \to 0{,}37$) mostra che la scelta della strategia di balancing
dipende dal contesto applicativo: se il costo di un bug sfuggito è prioritario (es.\
sicurezza), l'Undersampling è preferibile; se si vuole un bilancio complessivo,
l'Oversampling è la scelta migliore per RF. NaiveBayes è strutturalmente limitato dalla
violazione dell'assunzione di indipendenza condizionale: il bilanciamento non ne migliora
le performance perché il problema non è l'imbalance ma le correlazioni tra feature.

% TODO: completare le conclusioni con i risultati del run finale (SMOTE + wrapper FS +
% dataset aggiuntivi BOOKKEEPER/KAFKA). Aggiornare questa sezione al termine
% dell'esecuzione e integrare il confronto tra dataset multipli.

\subsection{Milestone 3}

% TODO: sintetizzare quante classi buggy avrebbero potuto essere prevenute e
% cosa implica questo risultato.

\subsection{Milestone 4}

% TODO: sintetizzare i risultati del refactoring automatico e le implicazioni pratiche.
